\documentclass[11pt]{article}
\usepackage[margin=1.15in]{geometry}
\usepackage{amsmath,amssymb,amsthm}
\usepackage{xcolor}
\usepackage[colorlinks=true,linkcolor=blue!55!black,citecolor=blue!55!black,urlcolor=blue!55!black]{hyperref}

\newtheorem{theorem}{Theorem}
\newtheorem{lemma}{Lemma}
\newtheorem{corollary}{Corollary}
\theoremstyle{definition}
\newtheorem{definition}{Definition}
\theoremstyle{remark}
\newtheorem{remark}{Remark}

\DeclareMathOperator{\lcm}{lcm}
\newcommand{\eps}{\varepsilon}

\title{\bfseries A Smoothness Covering Decomposition:\\
\large Covering $[1,n]$ by Divisibility with One Dedicated Element per Large Prime Power}
\author{Eric Schultz \and Claude (Anthropic)}
\date{July 2026}

\begin{document}
\maketitle

\begin{abstract}
We give an explicit, elementary construction of a set $D$ of positive integers
that covers $[1,n]$ by divisibility --- every $i\le n$ divides some element of
$D$ --- using a single ``smooth'' element together with one dedicated element
per prime power in $(\sqrt n,\,n]$. The construction has size
$|D|=\pi(n)\,(1+o(1))$ and its elements are of magnitude $\exp\!\big(O(\sqrt n)\big)$.
The proof rests on a single observation: no integer $i\le n$ has two coprime
maximal prime-power divisors both exceeding $\sqrt n$. The result isolates the
prime powers above $\sqrt n$ as the sole carriers of covering cost: every
$\sqrt n$-smooth number and every integer with a small cofactor is absorbed for
free by one shared element.
\end{abstract}

\section{Setup}

\begin{definition}
For a finite set $D\subset\mathbb Z_{>0}$, we say \emph{$D$ covers $[1,n]$ by
divisibility} if for every integer $i$ with $1\le i\le n$ there exists $d\in D$
with $i\mid d$.
\end{definition}

Throughout, $n\ge 2$ is an integer, $t=\lfloor\sqrt n\rfloor$, and
\[
  L \;=\; \lcm(1,2,\dots,t).
\]
For an integer $i\ge 1$ we call $p^{a}$ a \emph{maximal prime power of $i$} if
$p$ is prime, $a\ge 1$, and $p^{a}\parallel i$ (that is, $p^{a}\mid i$ but
$p^{a+1}\nmid i$). The maximal prime powers of $i$ are pairwise coprime and
their product equals $i$.

\section{The decomposition}

\begin{theorem}[Smoothness Covering Decomposition]\label{thm:main}
Let $n\ge 2$, $t=\lfloor\sqrt n\rfloor$, and $L=\lcm(1,\dots,t)$. Let
\[
  Q \;=\; \{\, q=p^{a} \;:\; p \text{ prime},\ a\ge 1,\ \sqrt n < q \le n \,\}
\]
be the set of prime powers in $(\sqrt n,\,n]$. Then
\[
  D \;=\; \{L\}\;\cup\;\{\,L\cdot q \;:\; q\in Q\,\}
\]
covers $[1,n]$ by divisibility, and
\[
  |D| \;=\; 1+|Q| \;=\; \pi(n)\,\bigl(1+o(1)\bigr).
\]
\end{theorem}

The proof uses the following elementary lemma.

\begin{lemma}\label{lem:atmostone}
Every integer $i$ with $1\le i\le n$ has at most one maximal prime power
exceeding $\sqrt n$.
\end{lemma}

\begin{proof}
Suppose $q_1=p_1^{a_1}$ and $q_2=p_2^{a_2}$ are two distinct maximal prime
powers of $i$, both greater than $\sqrt n$. Distinct maximal prime powers are
coprime, so $q_1 q_2\mid i$, whence $i\ge q_1 q_2 > \sqrt n\cdot\sqrt n = n$,
contradicting $i\le n$.
\end{proof}

\begin{proof}[Proof of Theorem~\ref{thm:main}]
Fix $i$ with $1\le i\le n$. By Lemma~\ref{lem:atmostone}, either all maximal
prime powers of $i$ are at most $\sqrt n$, or exactly one exceeds $\sqrt n$.

\medskip
\noindent\textbf{Case 1: every maximal prime power of $i$ is $\le\sqrt n$.}
Each maximal prime power $p^{a}\parallel i$ satisfies $p^{a}\le\sqrt n\le t$,\
hence $p^{a}\mid L$ because $L=\lcm(1,\dots,t)$ is divisible by every integer up
to $t$. The maximal prime powers of $i$ are pairwise coprime and their product
is $i$, so their product divides $L$; that is, $i\mid L$. Thus $i$ divides the
element $L\in D$.

\medskip
\noindent\textbf{Case 2: exactly one maximal prime power $q=p^{a}$ of $i$
satisfies $q>\sqrt n$.}
Then $q\in Q$ (as $q\mid i\le n$ gives $q\le n$), and $L\cdot q\in D$. Write
$i=q\,m$ with $m=i/q$. Since $q>\sqrt n$,
\[
  m=\frac{i}{q}\le\frac{n}{q}<\frac{n}{\sqrt n}=\sqrt n\le t,
\]
so $m\le t$ and therefore $m\mid L$. Moreover $q=p^{a}$ is the \emph{full} power
of $p$ dividing $i$, so $p\nmid m$, giving $\gcd(m,q)=1$. From $m\mid L$ and
$q\mid q$ with $\gcd(m,q)=1$ we obtain $m\,q\mid L\,q$. Since $i=m\,q$, it
follows that $i\mid L\cdot q$.

\medskip
In either case $i$ divides an element of $D$, so $D$ covers $[1,n]$.

For the count, $|D|=1+|Q|$ by construction. The elements of $Q$ are the prime
powers in $(\sqrt n,n]$. Prime powers $p^{a}$ with $a\ge 2$ and $p^a\le n$
number $O(\sqrt n)$ in total, while the primes in $(\sqrt n,n]$ number
$\pi(n)-\pi(\sqrt n)=\pi(n)(1+o(1))$. Hence $|Q|=\pi(n)(1+o(1))$ and
$|D|=\pi(n)(1+o(1))$.
\end{proof}

\section{Remarks}

\begin{remark}[Magnitude]
Since $\log L=\psi(t)=\psi(\lfloor\sqrt n\rfloor)\sim\sqrt n$ by the prime number
theorem (where $\psi$ is the second Chebyshev function), every element of $D$ has
magnitude
\[
  \max_{d\in D} d \;=\; L\cdot\max Q \;\le\; L\cdot n \;=\; \exp\!\big(\sqrt n\,(1+o(1))\big).
\]
\end{remark}

\begin{remark}[What the decomposition isolates]
The single element $L$ absorbs, for free, every $\sqrt n$-smooth integer and,
via Case~2, every integer of the form $m\,q$ with $m\le\sqrt n$ and $q$ a prime
power above $\sqrt n$. Consequently the only integers requiring their own
dedicated element are handled through the $\sim\pi(n)$ prime powers in
$(\sqrt n,n]$. In this sense the prime powers above $\sqrt n$ are the sole
carriers of covering cost in the construction.
\end{remark}

\begin{remark}[Verification]
The case analysis in the proof was checked exhaustively for every $i\le n$ and
all $n$ with $2\le n\le 3000$: in each instance the predicted covering element
($L$ in Case~1, $L\cdot q$ in Case~2) was confirmed to be divisible by $i$, the
``at most one large maximal prime power'' lemma held, and the full set $D$ was
confirmed to cover $[1,n]$.
\end{remark}

\end{document}
